\noindent\textbf{Elección del procedimiento.} El modelo lineal da $T_0=2\pi\sqrt{L/g}$. A amplitud de 90 grados, el parámetro de la integral elíptica es $m=\sin^2(\theta_0/2)=1/2$ y el periodo exacto del modelo conservativo es $T=4K(m)\sqrt{L/g}$. Se utiliza el valor de $K$ suministrado.

\begin{enumerate}

\item Calcular la referencia de pequeñas oscilaciones.

\[\fitmath{T_{0}=2\cdot \pi\cdot \sqrt{\frac{L}{g}}}\]

\[\fitmath{T_{0}=2\cdot \pi\cdot \sqrt{\frac{(1)}{(9{,}81)}}\approx 2{,}006067\,\mathrm{s}}\]

\item Introducir la integral elíptica del enunciado.

\[\fitmath{\mathrm{Texact}=4\cdot K\cdot \sqrt{\frac{L}{g}}}\]

\[\fitmath{\mathrm{Texact}=4\cdot (1{,}8540747)\cdot \sqrt{\frac{(1)}{(9{,}81)}}\approx 2{,}367842\,\mathrm{s}}\]

\item La fracción de exceso se multiplica por cien para dar porcentaje.

\[\fitmath{\mathrm{excess}=\frac{\mathrm{Texact}}{T_{0}}-1}\]

\[\fitmath{\mathrm{excess}=\frac{(2{,}367841977)}{(2{,}006066681)}-1\approx 0{,}1803406\,}\]

\end{enumerate}
